\documentclass[14pt,a4paper]{extarticle}
\usepackage{fontspec}
\usepackage{polyglossia}
\setmainlanguage{russian}
\setmainfont{Times New Roman}
\setmonofont{Courier New}
\usepackage{setspace}
\onehalfspacing
\usepackage{geometry}
\geometry{left=30mm, right=15mm, top=20mm, bottom=20mm}
\usepackage{indentfirst}
\setlength{\parindent}{1.25cm}
\setlength{\parskip}{0pt}
\usepackage{xurl}
\urlstyle{same}
\usepackage{hyperref}
\usepackage{booktabs}
\usepackage{ragged2e}
\usepackage{graphicx}
\usepackage{float}
\usepackage{pdflscape}
\usepackage{listings}
\usepackage{xcolor}
\usepackage{array}
\usepackage{amsmath} % <-- ИСПРАВЛЕНО: Подключен ДО wasysym
\usepackage{wasysym}
\usepackage{amsmath,amssymb,amsfonts}
\usepackage{geometry}
\usepackage{hyperref}
\usepackage{graphicx}
\usepackage{booktabs}
\usepackage{array}

\geometry{left=2cm,right=2cm,top=2cm,bottom=2cm}

\begin{document}

\part{Ответы на вопросы к коллоквиуму №1}

\section{Часть I. Элементарная теория погрешностей}

\subsection{1. Источники и классификация погрешностей. Абсолютная и относительная погрешности. Значащие цифры в десятичной записи числа. Округление. Верные значащие цифры. Правила записи приближённых чисел.}

\subsubsection{Источники и классификация погрешностей}

Погрешности обусловлены следующими причинами:

\begin{enumerate}
    \item \textbf{Приближённость математических моделей} --- любая модель является приближением реальности.
    \item \textbf{Погрешности исходных данных} --- неточности измерений.
    \item \textbf{Приближённость методов численного решения} --- численные методы сами по себе дают погрешности.
    \item \textbf{Потери точности при машинном представлении чисел} и арифметических операциях.
\end{enumerate}

Погрешности делятся на:
\begin{itemize}
    \item \textbf{Неустранимые} (причины 1 и 2) --- снижаются выбором более адекватной модели и более совершенных методов измерений.
    \item \textbf{Устранимые} (причины 3 и 4) --- уменьшаются совершенствованием методов, вычислительной техники и организации счёта.
\end{itemize}

\subsubsection{Абсолютная и относительная погрешности}

Пусть $x$ --- точное значение, $x^*$ --- приближённое значение.

\textbf{Погрешность:} $\Delta = x - x^*$

\textbf{Абсолютная погрешность:} $\Delta(x^*) = |x - x^*|$

Вычислить невозможно (точное значение неизвестно), поэтому используют \textbf{верхнюю оценку} $\overline{\Delta}(x^*)$, для которой заведомо известно: $\Delta(x^*) \leq \overline{\Delta}(x^*)$. Запись: $x = x^* \pm \overline{\Delta}(x^*)$.

\textbf{Относительная погрешность:}
\[\delta(x^*) = \frac{|x - x^*|}{|x^*|} = \frac{\Delta(x^*)}{|x^*|}\]

Также используют верхнюю оценку $\overline{\delta}(x^*)$. Относительную погрешность обычно дают в процентах.

Пример: при измерении длины спортзала и расстояния до метро абсолютная погрешность в обоих случаях может быть 1 м, но относительная погрешность длины спортзала --- 1--2\,\%, а расстояния до метро --- не более 0,01--0,1\,\%.

\subsubsection{Значащие цифры}

Приближённое число записывается в виде:
\[a = \pm m \cdot 10^r,\]
где $m = 0,\alpha_1\alpha_2\ldots\alpha_s$ --- мантисса, $r$ --- порядок.

\textbf{Значащие цифры} --- все цифры мантиссы, начиная с первой ненулевой слева. Остальные нули --- незначащие.

Примеры: $0{,}00\underline{1570}$, $\underline{25{,}030}$, $0{,}0\underline{70}$, $\underline{4{,}100000}$. Без незначащих нулей: $0{,}001570$, $25{,}03$, $0{,}070$, $4{,}100000$.

Замечания:
\begin{itemize}
    \item Значащие нули справа убирать нельзя: $1{,}50 \neq 1{,}500$ (разная точность).
    \item Для нуля: единственная значащая цифра --- последний нуль справа ($0 = 0{,}000$, значащая --- последний 0).
\end{itemize}

\subsubsection{Округление}

\textbf{Округление усечением} --- отбрасываются все значащие цифры правее нужной. Абсолютная погрешность не превышает единицы разряда, до которого округляли.

Примеры: $35{,}274$ до сотых $= 35{,}27$, до тысячных $= 35{,}274$; $2853$ до тысяч $= 2800$ (не 28, т.к. порядок сохраняется).

\textbf{Округление по дополнению} (применяется по умолчанию):
\begin{enumerate}
    \item Если первая отбрасываемая цифра $< 5$ --- оставляемые цифры не меняются.
    \item Если первая отбрасываемая цифра $> 5$, или $= 5$ и среди остальных есть ненулевые --- последняя сохраняемая цифра увеличивается на 1.
    \item Если первая отбрасываемая цифра $= 5$ и все остальные --- нули --- последняя сохраняемая цифра не меняется, если она чётная, и увеличивается на 1, если нечётная.
\end{enumerate}

Примеры: $0{,}1574$ до сотых $= 0{,}16$; $0{,}1575$ до тысячных $= 0{,}158$; $124{,}5$ до тысяч $= 1200$ (5 --- последняя, чётная 4 сохраняется).

Абсолютная погрешность не превосходит \textbf{половины единицы} разряда последней оставляемой цифры.

\subsubsection{Верные значащие цифры}

\textbf{Верная в широком смысле} --- значащая цифра, для которой абсолютная погрешность не превосходит \textbf{единицы} разряда, в котором она находится. Связана с округлением усечением.

\textbf{Верная в узком смысле} --- значащая цифра, для которой абсолютная погрешность не превосходит \textbf{половины единицы} разряда, в котором она находится. (По умолчанию.) Связана с округлением по дополнению.

Верная в узком смысле всегда верна и в широком, но не наоборот.

\textbf{Правило выявления верных цифр:}
\begin{enumerate}
    \item Абсолютная погрешность округляется с избытком до одной значащей цифры.
    \item Если эта цифра $\leq 5$ --- все разряды левее неё верные; если $> 5$ --- соседний слева разряд сомнительный, все левее --- верные.
\end{enumerate}

Примеры:
\begin{itemize}
    \item $x = 35{,}274$, $\overline{\Delta} = 0{,}03$ $\rightarrow$ округляем $\overline{\Delta}$ до одной значащей: $0{,}03 \to 0{,}03$; 3 $\leq$ 5 $\rightarrow$ верные цифры 3 и 5.
    \item $x = 284{,}5$, $\overline{\Delta} = 0{,}47$ $\rightarrow$ округляем: $0{,}47 \to 0{,}5$; 5 $\leq$ 5 $\rightarrow$ верная 2, сомнительна 8.
    \item $x = 1{,}357$, $\overline{\Delta} = 0{,}2$ $\rightarrow$ верной будет только 1.
\end{itemize}

\textbf{Обратная задача:} если $x^* = 2{,}740$ и все цифры верные, то $\overline{\Delta} = \frac{1}{2} \cdot 0{,}001 = 0{,}0005$. Если $x^* = 0{,}0781$ и все верные, то $\overline{\Delta} = \frac{1}{2} \cdot 0{,}0001 = 0{,}00005$.

При округлении до верных цифр нужно учитывать, что добавляется погрешность округления, и последняя верная цифра может стать сомнительной --- тогда округляют дальше справа налево.

\subsubsection{Погрешности представления чисел}

\textbf{Абсолютная погрешность представления:} $\Delta_M = |x - x^*_M|$, где $x^*_M$ --- машинное представление.

\textbf{Относительная погрешность представления:} $\delta_M = \dfrac{|x - x^*_M|}{|x^*_M|}$.

\textbf{Фиксированная запятая:} при длине разрядной сетки $p$ абсолютная погрешность лежит в диапазоне $[0; 2^{-p}]$, усреднённая $\approx 2^{-p-1}$. Относительная погрешность для малых чисел может быть очень значительной.

\textbf{Плавающая запятая (нормализованная форма):} мантисса в диапазоне $[0{,}5; 1)$, первый бит всегда 1. Относительная погрешность почти не зависит от величины числа. Поэтому математические вычисления над дробными числами проводят в форме с плавающей запятой.

\subsubsection{Правила записи приближённых чисел}

\begin{enumerate}
    \item В \textbf{промежуточных} результатах оставляют, кроме верных, \textbf{одну-две сомнительные} цифры.
    \item \textbf{Окончательный} результат округляют с сохранением \textbf{не более одной сомнительной} цифры.
\end{enumerate}

Причины не округлять строго до верных: оценки погрешностей завышены; желательно иметь запасную сомнительную цифру на случай погрешности округления.

\hrule\bigskip

\subsection{2. Погрешности арифметических операций}

\subsubsection{Теорема 1 (сумма и разность, абсолютная погрешность)}

\[\Delta(a^* \pm b^*) \leq \Delta(a^*) + \Delta(b^*)\]

Обобщение на $n$ величин: $\Delta\!\left(\sum a_i^*\right) \leq \sum \Delta(a_i^*)$.

\subsubsection{Теорема 2 (сумма и разность, относительная погрешность)}

Для ненулевых чисел одного знака:
\begin{itemize}
    \item \textbf{Сумма:} $\delta(a^* + b^*) \leq \max\{\delta(a^*), \delta(b^*)\}$
    \item \textbf{Разность:} $\delta(a^* - b^*) \leq \dfrac{\max\{\delta(a^*), \delta(b^*)\}}{q}$, где $q = \dfrac{|a^* - b^*|}{|a^*| + |b^*|}$
\end{itemize}

\textbf{Следствие:} при вычитании близких чисел ($a^* \approx b^*$) происходит \textbf{катастрофическая потеря точности} ($q \to 0$, погрешность $\to \infty$). Вместо $q$ используют приближение $q \approx \dfrac{|a^* - b^*|}{|a^*| + |b^*|}$.

\subsubsection{Теорема 3 (произведение)}

\[\delta(a^* \cdot b^*) \leq \delta(a^*) + \delta(b^*) + \delta(a^*)\delta(b^*)\]

Приближённо: $\delta(a^* \cdot b^*) \approx \delta(a^*) + \delta(b^*)$. Обобщение: $\delta(a_1^* \cdot \ldots \cdot a_n^*) \approx \sum \delta(a_i^*)$.

\subsubsection{Теорема 4 (частное)}

При $b^* \neq 0$, $\delta(b^*) < 1$:
\[\delta\!\left(\frac{a^*}{b^*}\right) \leq \frac{\delta(a^*) + \delta(b^*)}{1 - \delta(b^*)}\]

Приближённо: $\delta\!\left(\dfrac{a^*}{b^*}\right) \approx \delta(a^*) + \delta(b^*)$.

\subsubsection{Пример (подробное вычисление с соблюдением правил)}

Вычислить $u = \dfrac{x + y}{z - t}$, если $x = 3{,}45$ ($\overline{\Delta}_x = 0{,}005$), $y = 0{,}627$ ($\overline{\Delta}_y = 0{,}0005$), $z = 4{,}710$ ($\overline{\Delta}_z = 0{,}0005$), $t = 1{,}38$ ($\overline{\Delta}_t = 0{,}005$). Все цифры верные.

\begin{enumerate}
    \item Находим относительные погрешности:
    \begin{itemize}
        \item $\overline{\delta}_x = \dfrac{0{,}005}{3{,}45} \approx 0{,}15\%$
        \item $\overline{\delta}_y = \dfrac{0{,}0005}{0{,}627} \approx 0{,}08\%$
        \item $\overline{\delta}_z = \dfrac{0{,}0005}{4{,}710} \approx 0{,}01\%$
        \item $\overline{\delta}_t = \dfrac{0{,}005}{1{,}38} \approx 0{,}36\%$
    \end{itemize}
    \item Числитель $a = x + y = 4{,}077$. С одной запасной: $a = 4{,}08$.
    $\overline{\Delta}_a \leq 0{,}005 + 0{,}0005 = 0{,}0055 \approx 0{,}006$; $\overline{\delta}_a \leq \max(0{,}15\%, 0{,}08\%) = 0{,}15\%$.
    \item Знаменатель $b = z - t = 3{,}330$. С запасной: $b = 3{,}33$.
    $\overline{\Delta}_b \leq 0{,}0005 + 0{,}005 = 0{,}0055 \approx 0{,}006$; $\overline{\delta}_b = \dfrac{0{,}006}{3{,}33} \approx 0{,}18\%$.
    \item Результат: $u = \dfrac{4{,}08}{3{,}33} \approx 1{,}225$.
    $\overline{\delta}_u \approx 0{,}15\% + 0{,}18\% = 0{,}33\%$; $\overline{\Delta}_u = 1{,}225 \cdot 0{,}0033 \approx 0{,}004$.
    Ответ: $u = 1{,}23 \pm 0{,}004$, $\overline{\delta}_u \approx 0{,}33\%$ (с одной запасной сомнительной цифрой).
\end{enumerate}

\hrule\bigskip

\subsection{3. Погрешности функций. Линейная оценка дифференциальной погрешности.}

Пусть $f(x_1, \ldots, x_n)$ --- функция $n$ аргументов, $x^* = (x_1^*, \ldots, x_n^*)$ --- вектор приближённых значений с известными предельными погрешностями $\Delta(x_1^*), \ldots, \Delta(x_n^*)$.

Если $f$ непрерывно дифференцируема, то \textbf{линейная оценка дифференциальной погрешности:}

\textbf{Абсолютная погрешность:}
\[\Delta(f^*) \approx \sum_{i=1}^{n} \left|\frac{\partial f}{\partial x_i}(x^*)\right| \cdot \Delta(x_i^*)\]

\textbf{Относительная погрешность:}
\[\delta(f^*) \approx \sum_{i=1}^{n} \left|\frac{x_i^*}{f^*} \cdot \frac{\partial f}{\partial x_i}(x^*)\right| \cdot \delta(x_i^*) = \sum_{i=1}^{n} |p_i| \cdot \delta(x_i^*),\]
где $p_i = \dfrac{x_i^*}{f^*} \cdot \dfrac{\partial f}{\partial x_i}(x^*)$ --- показатели относительной погрешности.

Для функции одного аргумента $y = f(x)$:
\[\Delta(y^*) \approx |f'(x^*)| \cdot \Delta(x^*), \quad \delta(y^*) \approx \left|\frac{x^* f'(x^*)}{y^*}\right| \cdot \delta(x^*)\]

Числа обусловленности: $\vartheta_\Delta = |f'(x^*)|$, $\vartheta_\delta = \left|\dfrac{x^* f'(x^*)}{y^*}\right|$.

\subsubsection{Пример 1. Вычисление $u = \dfrac{\sqrt{x^3 - 1} - \ln y}{z}$}

Дано: $x = 3{,}020$ (все верные, $\overline{\Delta}_x = 0{,}0005$), $y = 0{,}4830$ (все верные, $\overline{\Delta}_y = 0{,}00005$).

Способ 1 (линейная оценка). Вычисляем $u^*$, частные производные $\dfrac{\partial u}{\partial x}$, $\dfrac{\partial u}{\partial y}$, $\dfrac{\partial u}{\partial z}$ в точке $(x^*, y^*, z^*)$, подставляем в формулу $\Delta(u^*) \approx \left|\dfrac{\partial u}{\partial x}\right|\overline{\Delta}_x + \left|\dfrac{\partial u}{\partial y}\right|\overline{\Delta}_y + \left|\dfrac{\partial u}{\partial z}\right|\overline{\Delta}_z$. Округляем результат до верных цифр с запасной.

Способ 2 (через относительную погрешность). Вычисляем $\overline{\delta}_x$, $\overline{\delta}_y$, $\overline{\delta}_z$, затем показатели $p_x, p_y, p_z$ и $\overline{\delta}_u \approx |p_x|\overline{\delta}_x + |p_y|\overline{\delta}_y + |p_z|\overline{\delta}_z$. Затем $\overline{\Delta}_u = |u^*| \cdot \overline{\delta}_u$.

Оба способа дают один результат: $u = \text{ответ с одной сомнительной цифрой}$.

\subsubsection{Пример 2. $f = \sum c_i x_i$ (линейная комбинация)}

Частные производные: $\dfrac{\partial f}{\partial x_i} = c_i$. Линейная оценка:
\[\Delta(f^*) \approx \sum |c_i| \cdot \Delta(x_i^*)\]

Для $f = x_1 + x_2$ или $f = x_1 - x_2$: $\Delta(f^*) \approx \Delta(x_1^*) + \Delta(x_2^*)$ --- получена оценка абсолютной погрешности суммы и разности.

Для $f = x_1 \cdot x_2$: $\delta(f^*) \approx \delta(x_1^*) + \delta(x_2^*)$ --- следствие из теоремы 3.

Для $f = \dfrac{x_1}{x_2}$: $\delta(f^*) \approx \delta(x_1^*) + \delta(x_2^*)$ --- следствие из теоремы 4.

\subsubsection{Пример 3. Обратная задача погрешности}

Дана $f(x, y) = \sqrt{x^2 + 1} - \ln y$. Вычислено при $x = 2$, $y = 3$. С какой точностью должны быть заданы аргументы, если $\Delta(f^*) \leq 0{,}001$?

\textbf{Принцип равного влияния:} вклад каждого аргумента в погрешность функции одинаков. Тогда:
\[\left|\frac{\partial f}{\partial x}\right| \Delta(x^*) = \left|\frac{\partial f}{\partial y}\right| \Delta(y^*) = \frac{\Delta(f^*)}{n}\]

Вычисляем частные производные в точке, находим допустимые $\Delta(x^*)$ и $\Delta(y^*)$.

\hrule\bigskip

\section{Часть II. Вычислительные задачи и методы}

\subsection{1. Особенности современных вычислительных задач. Этапы решения вычислительной задачи. Вычислительные задачи, их корректность и обусловленность. Классификация вычислительных методов.}

\subsubsection{Вычислительные задачи, корректность и обусловленность}

\textbf{Вычислительная задача} --- задание исходных числовых данных (входные данные) и формулировка требуемого числового результата (решения).

Решение называется \textbf{устойчивым} по отношению к малым возмущениям исходных данных, если последним соответствуют малые погрешности решения, не выше порядка погрешностей исходных данных. В простейшем скалярном случае: для любого $\varepsilon > 0$ существует $\delta(\varepsilon) > 0$, что для любого $x^*$ с $\Delta(x^*) < \delta$ соответствующее решение $y^*$ удовлетворяет $\Delta(y^*) < \varepsilon$. Решение непрерывно зависит от входных данных.

Вычислительная задача называется \textbf{корректной}, если:
\begin{enumerate}
    \item Для любого допустимого набора исходных данных существует \textbf{единственное} решение.
    \item Это решение \textbf{устойчиво} по отношению к малым возмущениям исходных данных.
\end{enumerate}

При невыполнении хотя бы одного условия задача называется \textbf{некорректной}.

\textbf{Число обусловленности} --- численная мера степени устойчивости решения, равная коэффициенту возможного роста погрешности решения по отношению к погрешностям исходных данных. В скалярном случае: если $\Delta y^* \leq \vartheta_\Delta \Delta x^*$, $\delta y^* \leq \vartheta_\delta \delta x^*$, то $\vartheta_\Delta$ --- абсолютное, $\vartheta_\delta$ --- относительное число обусловленности.

Для задачи вычисления функции $y = f(x)$:
\[\vartheta_\Delta = |f'(x^*)|, \quad \vartheta_\delta = \left|\frac{x^* f'(x^*)}{y^*}\right|\]

\begin{itemize}
    \item Если числа обусловленности \textbf{малы} (близки к нулю) --- задача \textbf{хорошо обусловлена}.
    \item Если \textbf{близки к единице или больше} --- задача \textbf{плохо обусловлена}.
\end{itemize}

\textbf{Пример 1.} Система:
\[\begin{cases} x_1 + x_2 = 0{,}5 \\ (1 + \varepsilon)x_1 + x_2 = 1{,}5 \end{cases}\]
Решение: $x_1 = 1/\varepsilon$, $x_2 = 0{,}5 - 1/\varepsilon$. При $\varepsilon = 0$ система не имеет решения, при $\varepsilon \neq 0$ --- имеет. Малое изменение параметра кардинально меняет решение --- задача плохо обусловлена.

\textbf{Пример 2.} Система:
\[\begin{cases} 1{,}1x_1 + x_2 = 1{,}1 \\ (1 + \varepsilon)x_1 + x_2 = 1 \end{cases}\]
Решение: $x_1 = \dfrac{1}{1 - 10\varepsilon}$, $x_2 = \dfrac{-11\varepsilon}{1 - 10\varepsilon}$. При $\varepsilon = 0{,}001$: $x_1 = 1{,}01$, $x_2 = -0{,}011$. При $\varepsilon = 0{,}002$: $x_1 = 1{,}02$, $x_2 = -0{,}022$. Абсолютное число обусловленности $x_2$: $\vartheta_\Delta = 11$ --- велико, задача плохо обусловлена.

Многие практические задачи некорректны. Существуют \textbf{методы регуляризации} некорректных задач (вклад А.Н. Тихонова). Один из подходов --- учёт априорной информации, сужающий множество решений до множества $M$, в котором решение устойчиво.

\subsubsection{Классификация вычислительных методов}

\begin{enumerate}
    \item \textbf{Методы эквивалентных преобразований} --- исходная задача заменяется эквивалентной, имеющей то же решение, но более простой. Например, решение $f(x) = 0$ заменяется поиском минимума $g(x) = (f(x))^2$.

    \item \textbf{Методы аппроксимации} --- исходная задача заменяется другой, решение которой близко к решению исходной. Оценка погрешности аппроксимации имеет большое значение. Например, вычисление интеграла заменяется интегрированием интерполяционного многочлена.

    \item \textbf{Прямые методы} --- точное решение получается за конечное число шагов. Примеры: формулы корней некоторых уравнений, метод Гаусса, метод прогонки для систем линейных уравнений.

    \item \textbf{Итерационные методы} --- строится последовательность приближений (итераций), каждая получается применением однотипного набора операций (итерационный шаг) к предыдущим.
    \begin{itemize}
        \item \textbf{Одношаговые} --- итерационный шаг использует одно приближение. Задаётся одно начальное приближение.
        \item \textbf{Многошаговые} --- итерационный шаг использует несколько приближений. Задаётся несколько начальных приближений (если невозможно --- сначала вычисляют одношаговым методом).
        \item Метод \textbf{сходится}, если итерационная последовательность сходится к точному решению. Множество начальных приближений, для которых метод сходится --- \textbf{область сходимости}.
        \item Процесс прекращается, когда достигается заданная точность (например, две последовательные итерации совпадают в пределах заданной величины).
    \end{itemize}
\end{enumerate}

Для итерационного метода необходимы:
\begin{itemize}
    \item Точное описание итерационного шага.
    \item Описание области сходимости или условия сходимости.
    \item Оценка погрешности текущей итерации для критерия останова.
\end{itemize}

\textbf{Пример:} извлечение корня $a = \sqrt{x}$. Итерационная схема Герона: $a_{k+1} = \frac{1}{2}\left(a_k + \dfrac{x}{a_k}\right)$. Для $x = 2$: 1 $\rightarrow$ 1,5 $\rightarrow$ 1,41667 $\rightarrow$ 1,41421. Сходится при любом положительном начальном приближении.

\subsubsection{Корректность алгоритмов}

\textbf{Вычислительный алгоритм} --- точное предписание действий над входными данными, задающее вычислительный процесс, направленный на преобразование произвольных входных данных в полностью определённый результат.

Алгоритм называется \textbf{корректным}, если:
\begin{enumerate}
    \item Он позволяет за конечное число операций преобразовать любое допустимое входное данное $x \in X$ в единственный результат $y \in Y$ ($X$ --- множество допустимых данных, $Y$ --- множество решений).
    \item Полученное решение \textbf{устойчиво} по отношению к малым возмущениям входных данных.
    \item Результат обладает \textbf{вычислительной устойчивостью}.
\end{enumerate}

\textbf{Устойчивость по входным данным} --- результат непрерывным образом зависит от входных данных при отсутствии вычислительной погрешности.

\textbf{Вычислительная устойчивость} --- стремление к нулю вычислительной погрешности (обусловленной приближённостью метода и машинной погрешностью) при стремлении к нулю машинной погрешности.

Алгоритм \textbf{вычислительно устойчив}, если результат вычислительно устойчив при любых допустимых входных данных. Алгоритм \textbf{устойчив по входным данным}, если результат устойчив по входным данным. Если алгоритм устойчив по входным данным и вычислительно устойчив --- он называется \textbf{устойчивым}. Алгоритм корректен, если он даёт единственное решение за конечное число шагов при любых допустимых данных и устойчив.

\textbf{Число обусловленности алгоритма:} если $\delta y^* \leq \vartheta_A \cdot \varepsilon_M$, где $\varepsilon_M$ --- относительная погрешность округления, то $\vartheta_A$ --- число обусловленности алгоритма. Если $\vartheta_A \approx 0$ --- алгоритм хорошо обусловлен, если $\vartheta_A \gg 1$ --- плохо обусловлен.

\hrule\bigskip

\section{Часть III. Аналитическое приближение табличных функций}

\subsection{1. Задача приближённого вычисления функции. Интерполяция. Задача полиномиальной интерполяции. Теорема о единственности решения. Погрешность полиномиальной интерполяции.}

\subsubsection{Задача приближённого вычисления функции}

Функция $y = f(x)$ задана таблицей $(x_i, y_i)$, $i = 0, \ldots, n$. Требуется приближённо вычислить значение в точке $x$, не совпадающей ни с одним из узлов.

Строится просто вычисляемая функция $y = g(x)$, значения которой в узлах совпадают с табличными. Вне узлов считается, что $f(x) \approx g(x)$. Это --- \textbf{интерполяция}.

\subsubsection{Интерполяция}

Функция $g$ называется \textbf{интерполирующей}, значения $x_0, \ldots, x_n$ --- \textbf{узлами интерполяции}. Условия $g(x_i) = y_i$ называются \textbf{условиями интерполяции}.

\subsubsection{Полиномиальная интерполяция}

В качестве класса интерполирующих функций берётся множество алгебраических многочленов степени $m$:
\[g(x) = P_m(x) = a_0 + a_1 x + \cdots + a_m x^m.\]

Задача: найти коэффициенты $a_0, \ldots, a_m$, удовлетворяющие условиям интерполяции $P_m(x_i) = y_i$.

\textbf{Теорема о единственности:} если $f$ задана таблицей с попарно различными узлами, то существует \textbf{единственный} интерполяционный многочлен $P_n$ степени $n$ (при $m = n$).

\textit{Доказательство:} коэффициенты находятся из системы $n+1$ линейных уравнений с $n+1$ неизвестным. Её определитель --- определитель Вандермонда, который не равен нулю при попарно различных узлах. Следовательно, система имеет единственное решение.

Замечания:
\begin{itemize}
    \item При $m > n$ --- бесконечное множество решений.
    \item При $m < n$ --- может вообще не иметь решений.
    \item Степень может быть меньше $n$, если старшие коэффициенты нулевые.
\end{itemize}

\subsubsection{Погрешность полиномиальной интерполяции}

\textbf{Теорема:} пусть функция $f$ имеет $n+1$ непрерывную производную на отрезке $[a; b]$, содержащем все узлы. Тогда для любого $x \in [a; b]$ существует $\zeta \in (a; b)$:
\[R_n(x) = f(x) - P_n(x) = \frac{f^{(n+1)}(\zeta)}{(n+1)!} Q_{n+1}(x),\]
где $Q_{n+1}(x) = (x - x_0)(x - x_1)\cdots(x - x_n)$.

\textbf{Оценочная функция:}
\[\Delta(P_n(x)) = |R_n(x)| \leq \frac{M_{n+1}}{(n+1)!} |Q_{n+1}(x)|,\]
где $M_{n+1} = \max\limits_{z \in [a; b]} |f^{(n+1)}(z)|$.

\textbf{Глобальная оценка:} $V(P_n) = \sup\limits_{x \in [a; b]} \Delta(P_n(x))$.

На практике вместо $M_{n+1}$ используют её верхнюю оценку $q_{n+1}$.

\textbf{Важно:} с ростом $n$ глобальная погрешность может неограниченно возрастать (\textbf{феномен Рунге}). Поэтому на практике интерполяционные многочлены степени выше третьей не применяются.

\hrule\bigskip

\subsection{2. Интерполяционный многочлен Лагранжа. Схема Эйткена.}

\subsubsection{Интерполяционный многочлен Лагранжа}

Формула:
\[L_n(x) = \sum_{i=0}^{n} y_i \cdot l_{n,i}(x),\]
где множители Лагранжа:
\[l_{n,i}(x) = \prod_{\substack{j=0 \\ j \neq i}}^{n} \frac{x - x_j}{x_i - x_j}.\]

Для $n = 1$ (два узла):
\[L_1(x) = y_0 \frac{x - x_1}{x_0 - x_1} + y_1 \frac{x - x_0}{x_1 - x_0}\]
--- уравнение прямой через две точки.

Для $n = 2$ (три узла):
\[L_2(x) = y_0 \frac{(x-x_1)(x-x_2)}{(x_0-x_1)(x_0-x_2)} + y_1 \frac{(x-x_0)(x-x_2)}{(x_1-x_0)(x_1-x_2)} + y_2 \frac{(x-x_0)(x-x_1)}{(x_2-x_0)(x_2-x_1)}\]
--- уравнение параболы через три точки.

\textbf{Недостаток:} при добавлении нового узла формулу надо пересчитывать заново.

\textbf{Пример.} Функция $f(t) = t^4 + 2t^3 - 7t^2 - 8t + 12$ задана 4 точками:

\begin{center}
\begin{tabular}{c|c|c|c|c}
\toprule
$x_i$ & $-1$ & $0$ & $1$ & $3$ \\
\midrule
$y_i$ & $12$ & $12$ & $0$ & $60$ \\
\bottomrule
\end{tabular}
\end{center}

\[L_3(x) = 12 \cdot \frac{(x-0)(x-1)(x-3)}{(-1-0)(-1-1)(-1-3)} + 12 \cdot \frac{(x+1)(x-1)(x-3)}{(0+1)(0-1)(0-3)} + 60 \cdot \frac{(x+1)(x-0)(x-1)}{(3+1)(3-0)(3-1)}\]

После упрощения: $L_3(x) = 5x^3 - 6x^2 - 11x + 12$. Исходная функция --- многочлен 4-й степени, а приближение --- 3-й степени. На отрезке $[1; 3]$ погрешность высока.

\subsubsection{Схема Эйткена}

Рекуррентная схема пошагового вычисления значения $L_n(x)$ в точке $x$.

Обозначим $L_k^{i, i+1, \ldots, i+k}(x)$ --- значение полинома Лагранжа степени $k$ по узлам $x_i, \ldots, x_{i+k}$ в точке $x$. Нулевой шаг: $L_0^i(x) = y_i$.

\textbf{Расчётная формула $k$-го шага:}
\[L_k^{i, \ldots, i+k}(x) = \frac{1}{x_{i+k} - x_i} \begin{vmatrix} L_{k-1}^{i, \ldots, i+k-1}(x) & x_i - x \\ L_{k-1}^{i+1, \ldots, i+k}(x) & x_{i+k} - x \end{vmatrix}\]
\[= \frac{L_{k-1}^{i, \ldots, i+k-1}(x)(x_{i+k} - x) - L_{k-1}^{i+1, \ldots, i+k}(x)(x_i - x)}{x_{i+k} - x_i}\]

\begin{itemize}
    \item На 1-м шаге: по парам соседних узлов строятся многочлены 1-й степени.
    \item На 2-м шаге: по тройкам --- многочлены 2-й степени.
    \item На последнем, $n$-м шаге: вычисляется $L_n(x)$.
\end{itemize}

\textbf{Доказательство} (индукция по $k$): база $k = 1$ даёт многочлен Лагранжа первой степени. Индукционный переход: формула даёт многочлен $(k+1)$-й степени, удовлетворяющий условиям интерполяции во всех узлах.

\textbf{Пример} (та же функция, $x = 0{,}6$). Таблица Эйткена:

\begin{center}
\begin{tabular}{c|c|c|c|c|c}
\toprule
$x_i$ & $y_i$ & $x_i - x$ & Шаг 1 & Шаг 2 & Шаг 3 \\
\midrule
$-1$ & $12$ & $-1{,}6$ & & & \\
$0$ & $12$ & $-0{,}6$ & $4{,}80$ & & \\
$1$ & $0$ & $0{,}4$ & $-7{,}20$ & $3{,}60$ & \\
$3$ & $60$ & $2{,}4$ & $36{,}00$ & $1{,}80$ & $4{,}32$ \\
\bottomrule
\end{tabular}
\end{center}

Результат: $L_3(0{,}6) = 4{,}32$. Проверка: $L_3(0{,}6) = 5(0{,}6)^3 - 6(0{,}6)^2 - 11(0{,}6) + 12 = 4{,}32$. \checkmark

\hrule\bigskip

\subsection{3. Разделённые разницы и их свойства. Интерполяционные многочлены Ньютона с разделёнными разностями.}

\subsubsection{Разделённые разности}

\textbf{Разделённая разность 1-го порядка:}
\[f(x_i; x_{i+1}) = \frac{y_{i+1} - y_i}{x_{i+1} - x_i} = \frac{f(x_{i+1}) - f(x_i)}{x_{i+1} - x_i}\]

\textbf{Разделённая разность $k$-го порядка} (рекуррентно):
\[f(x_i; \ldots; x_{i+k}) = \frac{f(x_{i+1}; \ldots; x_{i+k}) - f(x_i; \ldots; x_{i+k-1})}{x_{i+k} - x_i}\]

\textbf{Формула для непосредственного вычисления:}
\[f(x_i; \ldots; x_{i+k}) = \sum_{j=i}^{i+k} \frac{f(x_j)}{\prod\limits_{\substack{m=i \\ m \neq j}}^{i+k} (x_j - x_m)}\]

\textbf{Свойства:}
\begin{enumerate}
    \item Разделённые разности --- \textbf{симметричные функции} своих аргументов (не зависят от перестановок узлов).
    \item Являются дискретным аналогом производной.
\end{enumerate}

\subsubsection{Многочлен Ньютона с разделёнными разностями (первый вид, вперёд)}

\[P_n^{(1)}(x) = f(x_0) + f(x_0; x_1)(x - x_0) + f(x_0; x_1; x_2)(x - x_0)(x - x_1) + \cdots + f(x_0; \ldots; x_n)(x - x_0)\cdots(x - x_{n-1})\]
\[= f(x_0) + \sum_{i=1}^{n} f(x_0; \ldots; x_i)(x - x_0)\cdots(x - x_{i-1})\]

Разделённые разности берутся из \textbf{верхней диагонали} таблицы.

\subsubsection{Многочлен Ньютона с разделёнными разностями (второй вид, назад)}

\[P_n^{(2)}(x) = f(x_n) + f(x_{n-1}; x_n)(x - x_n) + f(x_{n-2}; x_{n-1}; x_n)(x - x_n)(x - x_{n-1}) + \cdots + f(x_0; \ldots; x_n)(x - x_n)\cdots(x - x_1)\]
\[= f(x_n) + \sum_{i=1}^{n} f(x_{n-i}; \ldots; x_n)(x - x_n)\cdots(x - x_{n-i+1})\]

Разделённые разности берутся из \textbf{нижней косой строки} таблицы.

\textbf{Преимущество перед Лагранжем:} при добавлении нового узла достаточно вычислить и прибавить одно новое слагаемое с разделённой разностью высшего порядка. Новый узел добавляется в конец таблицы (для первого вида) или в начало (для второго вида), остальные разности не меняются.

\textbf{Оценка погрешности (приближённая):}
\[\Delta(P_n(x)) \approx |f(x; x_0; \ldots; x_n) \cdot Q_{n+1}(x)|,\]
где $f(x; x_0; \ldots, x_n)$ --- разделённая разность $(n+1)$-го порядка с добавленной точкой $x$. Точку $x$ можно поставить в любое место таблицы (обычно до первого или после последнего узла --- достаточно рассчитать одну строку).

\begin{itemize}
    \item $P_n^{(1)}$ лучше для интерполяции в \textbf{начале} таблицы.
    \item $P_n^{(2)}$ лучше для интерполяции в \textbf{конце} таблицы.
    \item При неполных таблицах: первый вид применяется при интерполяции в начале (меньшая погрешность), второй --- в конце.
\end{itemize}

\textbf{Пример} (та же функция $f(t) = t^4 + 2t^3 - 7t^2 - 8t + 12$):

Таблица разделённых разностей:

\begin{center}
\begin{tabular}{c|c|c|c|c}
\toprule
$x_i$ & $f(x_i)$ & 1-й порядок & 2-й порядок & 3-й порядок \\
\midrule
$-1$ & $12$ & & & \\
& & $0$ & & \\
$0$ & $12$ & & $-6$ & \\
& & $-12$ & & $5$ \\
$1$ & $0$ & & $14$ & \\
& & $30$ & & \\
$3$ & $60$ & & & \\
\bottomrule
\end{tabular}
\end{center}

\textbf{Первый вид} (верхняя диагональ):
\[P_3^{(1)}(x) = 12 - 6(x + 1)(x - 0) + 5(x + 1)(x - 0)(x - 1) = 5x^3 - 6x^2 - 11x + 12\]

\textbf{Второй вид} (нижняя косая строка):
\[P_3^{(2)}(x) = 60 + 30(x - 3) + 14(x - 3)(x - 1) + 5(x - 3)(x - 1)(x - 0) = 5x^3 - 6x^2 - 11x + 12\]

Оба многочлена совпадают (теорема единственности).

\textbf{Добавление узла:} если построен $P_2^{(2)}$ по узлам $x_1, x_2, x_3$: $P_2^{(2)}(x) = 60 + 30(x-3) + 14(x-3)(x-1) = 14x^2 - 26x + 12$, то при добавлении $x_0$:
\[P_3^{(2)}(x) = P_2^{(2)}(x) + f(x_0; x_1; x_2; x_3)(x - 3)(x - 1)(x - 0) = 14x^2 - 26x + 12 + 5(x-3)(x-1)(x)\]
Получаем тот же результат.

\hrule\bigskip

\subsection{4. Конечные разности и их свойства. Интерполяционные многочлены Ньютона с конечными разностями.}

\subsubsection{Конечные разности}

Определяются для \textbf{равномерных} таблиц ($h = x_i - x_{i-1} = \text{const}$, $x_i = x_0 + ih$).

\textbf{Конечная разность 1-го порядка:}
\[\Delta y_i = y_{i+1} - y_i\]

\textbf{Конечная разность $k$-го порядка} (рекуррентно):
\[\Delta^k y_i = \Delta^{k-1} y_{i+1} - \Delta^{k-1} y_i\]

\textbf{Формула через биномиальные коэффициенты:}
\[\Delta^k y_i = \sum_{j=0}^{k} (-1)^j C_k^j \, y_{i+k-j}, \quad C_k^j = \frac{k!}{j!\,(k-j)!}\]

\textbf{Связь с разделёнными разностями:}
\[f(x_i; \ldots; x_{i+k}) = \frac{\Delta^k y_i}{k! \, h^k}\]

\subsubsection{Многочлен Ньютона с конечными разностями (первый вид, вперёд)}

\[P_n^{(1)}(x) = f(x_0) + \frac{\Delta y_0}{1! \, h}(x - x_0) + \frac{\Delta^2 y_0}{2! \, h^2}(x - x_0)(x - x_1) + \cdots + \frac{\Delta^n y_0}{n! \, h^n}(x - x_0)\cdots(x - x_{n-1})\]

С заменой $q = \frac{x - x_0}{h}$ (т.е. $x = x_0 + qh$):
\[P_n^{(1)}(x_0 + qh) = f(x_0) + \frac{\Delta y_0}{1!} q + \frac{\Delta^2 y_0}{2!} q(q-1) + \cdots + \frac{\Delta^n y_0}{n!} q(q-1)\cdots(q-n+1)\]

\subsubsection{Многочлен Ньютона с конечными разностями (второй вид, назад)}

\[P_n^{(2)}(x) = f(x_n) + \frac{\Delta y_{n-1}}{1! \, h}(x - x_n) + \frac{\Delta^2 y_{n-2}}{2! \, h^2}(x - x_n)(x - x_{n-1}) + \cdots + \frac{\Delta^n y_0}{n! \, h^n}(x - x_n)\cdots(x - x_1)\]

С заменой $q = \frac{x - x_n}{h}$:
\[P_n^{(2)}(x_n + qh) = f(x_n) + \frac{\Delta y_{n-1}}{1!} q + \frac{\Delta^2 y_{n-2}}{2!} q(q+1) + \cdots + \frac{\Delta^n y_0}{n!} q(q+1)\cdots(q+n-1)\]

\textbf{Оценочные функции погрешности:}
\[\Delta(P_n^{(1)}(x_0 + qh)) = \frac{M_{n+1}}{(n+1)!} h^{n+1} |q(q-1)\cdots(q-n)|\]
\[\Delta(P_n^{(2)}(x_n + qh)) = \frac{M_{n+1}}{(n+1)!} h^{n+1} |q(q+1)\cdots(q+n)|\]

\hrule\bigskip

\subsection{5. Интерполяционные многочлены Гаусса, Стирлинга и Бесселя.}

Применяются для \textbf{интерполяции в середине} равномерной таблицы. Узлы: $a, a \pm h, a \pm 2h, \ldots$, точка $x$: $a < x < a + h$.

Все выводятся из многочлена Ньютона с конечными разностями.

\subsubsection{Многочлен Гаусса (вперёд)}

Строится по узлам $x_0 = a, x_1 = a + h, x_{-1} = a - h, x_2 = a + 2h, \ldots$ с использованием разностей из таблицы, расположенных ``зигзагом'' вверх-вниз от центральной строки.

\[G_n(x) = y_0 + \frac{\Delta y_0}{1!} q + \frac{\Delta^2 y_{-1}}{2!} q(q-1) + \frac{\Delta^3 y_{-1}}{3!} q(q^2 - 1) + \cdots\]

\subsubsection{Многочлен Стирлинга}

Среднее арифметическое прямого и обратного многочленов Гаусса:
\[S_n(x) = \frac{G_n^{\text{вперёд}} + G_n^{\text{назад}}}{2}\]

\[S_n(x) = y_0 + \frac{\Delta y_{-1} + \Delta y_0}{2 \cdot 1!} q + \frac{\Delta^2 y_{-1}}{2!} q^2 + \frac{\Delta^3 y_{-2} + \Delta^3 y_{-1}}{2 \cdot 3!} q(q^2 - 1) + \cdots\]

Использует \textbf{средние арифметические} конечных разностей нечётных порядков.

\subsubsection{Многочлен Бесселя}

Применяется при $0.25 \leq q \leq 0.75$ (середина межузелья). Строится по узлам, симметричным относительно середины.

\[B_n(x) = \frac{y_0 + y_1}{2} + \left(q - \frac{1}{2}\right) \Delta y_0 + \frac{q(q-1)}{2!} \frac{\Delta^2 y_{-1} + \Delta^2 y_0}{2} + \cdots\]

\textbf{Практическое правило:} для интерполяции в начале таблицы --- $P_n^{(1)}$ (Ньютон вперёд) или Гаусс вперёд; в конце --- $P_n^{(2)}$ (Ньютон назад); в середине --- Стирлинг или Бессель.

\hrule\bigskip

\subsection{6. Задача аппроксимации табличной функции. Аппроксимация в классе обобщённых многочленов методом наименьших квадратов. Полиномиальная аппроксимация.}

\subsubsection{Задача аппроксимации}

В отличие от интерполяции, аппроксимирующая функция \textbf{приближённо} повторяет табличные данные, отслеживая общую тенденцию. Условия:
\[g(x_i) \approx y_i, \quad i = 1, \ldots, n.\]

Это необходимо, когда:
\begin{itemize}
    \item Таблицы большие, интерполяционные многочлены имеют высокую степень и сложны.
    \item Данные содержат ошибки --- желательно их \textbf{сглаживать}, а не точно отслеживать.
\end{itemize}

\subsubsection{Метод наименьших квадратов (МНК)}

Строится функция $g \in \mathbb{G}$, минимизирующая \textbf{среднеквадратическое отклонение}:
\[\delta^2(g) = \sqrt{\frac{1}{n} \sum_{i=1}^{n} (y_i - g(x_i))^2} \to \min\]

Пусть $\mathbb{G}$ --- множество \textbf{обобщённых многочленов} степени $m$ на системе линейно независимых базисных функций $\{\varphi_0, \varphi_1, \ldots, \varphi_m\}$:
\[\Phi_m(x) = \sum_{j=0}^{m} a_j \varphi_j(x)\]

Минимизируется:
\[S(a_0; \ldots; a_m) = \sum_{i=1}^{n} \left(y_i - \sum_{j=0}^{m} a_j \varphi_j(x_i)\right)^2 \to \min\]

Взяв частные производные и приравняв к нулю, получаем \textbf{нормальную систему}:
\[\sum_{j=0}^{m} a_j \sum_{i=1}^{n} \varphi_j(x_i) \varphi_k(x_i) = \sum_{i=1}^{n} y_i \varphi_k(x_i), \quad k = 0, \ldots, m\]

В матричном виде: $\Gamma a = b$, где $\Gamma = P^T P$ --- \textbf{матрица Грама}.

Если матрица Грама невырожденная (что верно при попарно несовпадающих узлах), система имеет единственное решение, которое является точкой минимума $S$.

\subsubsection{Полиномиальная аппроксимация}

Частный случай: $\varphi_j(x) = x^j$ (степенная система). Тогда $\Phi_m(x) = a_0 + a_1 x + \cdots + a_m x^m = P_m(x)$.

Нормальная система:
\[\sum_{j=0}^{m} a_j \sum_{i=1}^{n} x_i^{j+k} = \sum_{i=1}^{n} y_i x_i^k, \quad k = 0, \ldots, m\]

Для \textbf{линейной аппроксимации} ($m = 1$):
\[\begin{cases} a_0 n + a_1 \sum x_i = \sum y_i \\ a_0 \sum x_i + a_1 \sum x_i^2 = \sum y_i x_i \end{cases}\]

Для \textbf{квадратичной аппроксимации} ($m = 2$):
\[\begin{cases} a_0 n + a_1 \sum x_i + a_2 \sum x_i^2 = \sum y_i \\ a_0 \sum x_i + a_1 \sum x_i^2 + a_2 \sum x_i^3 = \sum y_i x_i \\ a_0 \sum x_i^2 + a_1 \sum x_i^3 + a_2 \sum x_i^4 = \sum y_i x_i^2 \end{cases}\]

Наилучшее среднеквадратическое отклонение:
\[\delta^2(P_2) = \sqrt{\frac{1}{n}\sum_{i=1}^{n} (y_i - a_0 - a_1 x_i - a_2 x_i^2)^2}\]

Геометрически линейная и квадратичная аппроксимация --- это наилучшее сглаживание набора точек прямой и параболой соответственно.

Замечание: при $m = n - 1$ аппроксимирующий полином совпадает с интерполяционным ($\delta^2 = 0$). Задача интерполяции --- частный случай аппроксимации.

\textbf{Пример.} Функция $y = \sin x$ аппроксимируется по 15 точкам отрезка $[1; 3{,}8]$ с шагом $h = 0{,}2$:
\begin{itemize}
    \item $P_1(x) = 1{,}856 - 0{,}586x$, $\delta^2(P_1) = 0{,}198$
    \item $P_2(x) = 0{,}414 + 0{,}795x - 0{,}288x^2$, $\delta^2(P_2) = 0{,}055$
\end{itemize}

Парабола даёт лучшее приближение (меньшее среднеквадратическое отклонение), чем прямая. Вычисления проводятся через матрицу Грама $A = P^T P$ и вектор правой части $b = P^T y$.

\hrule\bigskip

\subsection{7. Неполиномиальная аппроксимация.}

МНК можно применять не только для полиномов. Если есть предположения о виде функции, задаётся семейство функций с параметрами, которые подбираются из условия минимума среднеквадратического отклонения.

Часто используемые семейства:
\[y = a + b \ln x, \quad y = a + b \lg x, \quad y = ax^b, \quad y = ae^{bx}\]
\[y = b + \frac{a}{x}, \quad y = \frac{1}{ax + b}, \quad y = \frac{x}{ax + b}\]

\textbf{Пример 1:} $y = a + \dfrac{b}{x^2}$. Минимизируем $S(a, b) = \sum \left(y_i - a - \dfrac{b}{x_i^2}\right)^2$. Система линейная:
\[\begin{cases} an + b \sum \dfrac{1}{x_i^2} = \sum y_i \\ a \sum \dfrac{1}{x_i^2} + b \sum \dfrac{1}{x_i^4} = \sum \dfrac{y_i}{x_i^2} \end{cases}\]

\textbf{Пример 2:} $y = ae^{bx}$. Система \textbf{нелинейная}:
\[\begin{cases} \sum (y_i - ae^{bx_i}) e^{bx_i} = 0 \\ \sum (y_i - ae^{bx_i}) x_i e^{bx_i} = 0 \end{cases}\]
Решается приближёнными итерационными методами.

\textbf{Сложность:} для неполиномиальной аппроксимации МНК не всегда обоснован --- решение нормальной системы может не существовать, а если существует, то не обязано быть точкой минимума. В каждом случае надо устанавливать существование решения и проверять на минимум.

\hrule\bigskip

\subsection{8. Многочлены Чебышёва, их свойства.}

\subsubsection{Определения}

\textbf{Рекуррентное:}
\[T_0(x) = 1, \quad T_1(x) = x, \quad T_n(x) = 2x T_{n-1}(x) - T_{n-2}(x), \quad n \geq 2\]

Примеры:
\[T_2(x) = 2x^2 - 1, \quad T_3(x) = 4x^3 - 3x, \quad T_4(x) = 8x^4 - 8x^2 + 1\]

\textbf{Тригонометрическое} (при $|x| \leq 1$):
\[T_n(x) = \cos(n \arccos x)\]

\textbf{Аналитическое} (при $|x| > 1$):
\[T_n(x) = \frac{1}{2}\left(\left(x + \sqrt{x^2 - 1}\right)^n + \left(x - \sqrt{x^2 - 1}\right)^n\right)\]

\subsubsection{Свойства}

\begin{enumerate}
    \item \textbf{Старший коэффициент} $T_n$ равен $2^{n-1}$: $T_n(x) = 2^{n-1}x^n + \cdots$

    \item \textbf{Чётность:} $T_{2n}$ --- чётные функции, $T_{2n+1}$ --- нечётные.

    \item \textbf{Корни} на $[-1; 1]$:
    \[x_m = \cos\frac{\pi(2m - 1)}{2n}, \quad m = 1, \ldots, n\]

    \item \textbf{Точки экстремума} на $[-1; 1]$:
    \[\tilde{x}_m = \cos\frac{\pi m}{n}, \quad m = 0, \ldots, n\]
    Причём $T_n(\tilde{x}_m) = (-1)^m$, т.е. $|T_n(x)| \leq 1$ при $|x| \leq 1$.

    \item \textbf{Наименьшее уклонение от нуля:} нормированный многочлен Чебышёва $\overline{T}_n(x) = 2^{1-n} T_n(x) = x^n + \cdots$ наименее уклоняется от нуля среди всех нормированных многочленов степени $n$ на $[-1; 1]$:
    \[\max_{x \in [-1; 1]} |P_n(x)| \geq \max_{x \in [-1; 1]} |\overline{T}_n(x)| = 2^{1-n}\]

    \item Для произвольного отрезка $[a; b]$ (замена $x = \frac{2z - (a+b)}{b-a}$):
    \[\overline{T}_n^{[a; b]}(x) = (b-a)^n 2^{1-2n} T_n\!\left(\frac{2x - (a + b)}{b - a}\right)\]
    \[\max_{x \in [a; b]} |P_n(x)| \geq (b-a)^n 2^{1-2n}\]

    Корни:
    \[x_m = \frac{a + b}{2} + \frac{b - a}{2} \cos\frac{\pi(2m - 1)}{2n}, \quad m = 1, \ldots, n\]
\end{enumerate}

\subsubsection{Построение многочлена Чебышёва тремя способами:}
\begin{enumerate}
    \item По рекуррентной формуле $T_n = 2x T_{n-1} - T_{n-2}$.
    \item По тригонометрической формуле $T_n(x) = \cos(n \arccos x)$ с раскрытием через $\cos n\theta$.
    \item По аналитической формуле через корни характеристического уравнения.
\end{enumerate}

\hrule\bigskip

\subsection{9. Задача минимизации оценки остаточного члена интерполяции. Теоремы Фабера и Марцинкевича. Решение с помощью многочленов Чебышёва.}

\subsubsection{Феномен Рунге}

При использовании равномерной сетки с большим числом узлов глобальная погрешность может неограниченно возрастать:
\[\lim_{n \to \infty} \|f - P_n\| = \infty\]

\textbf{Функция Рунге:} $f(x) = \dfrac{1}{1 + 25x^2}$ на $[-1; 1]$. Бесконечно дифференцируема, но производные быстро растут:

\begin{center}
\begin{tabular}{c|c|c|c|c|c|c}
\toprule
$n$ & 1 & 2 & 3 & 4 & 5 & 6 \\
\midrule
$|f^{(n)}(1)|$ & 0,08 & 0,22 & 0,80 & 3,64 & 19,77 & 250 \\
\bottomrule
\end{tabular}
\end{center}

Оценка погрешности $V(P_n)$ увеличивается с ростом $n$. На графиках интерполяционных многочленов $P_n$ ($n = 3, 5, 7, 12$) видны значительные отклонения вблизи концов интервала, увеличивающиеся с ростом $n$.

\subsubsection{Теорема Фабера}

Какова бы ни была стратегия выбора узлов интерполяции, найдётся непрерывная на $[a; b]$ функция $f$, для которой
\[\lim_{n \to \infty} \|f - P_n\| = \infty\]

Это отрицает возможность решения при \textbf{фиксированной} стратегии построения сетки (например, равномерной).

\subsubsection{Теорема Марцинкевича}

Для \textbf{любой} непрерывной на $[a; b]$ функции $f$ существует такая последовательность интерполяционных сеток, что
\[\lim_{n \to \infty} \|f - P_n\| = 0\]

Решение существует для любой непрерывной функции --- нужно правильно подобрать узлы.

\subsubsection{Решение через узлы Чебышёва}

Оценка погрешности:
\[\|f - P_n\| \leq \frac{M_{n+1}}{(n+1)!} \|Q_{n+1}\|\]

Многочлен $Q_{n+1}(x) = (x - x_0)(x - x_1)\cdots(x - x_n)$ --- нормированный. Минимум $\|Q_{n+1}\|$ достигается, когда $Q_{n+1}$ совпадает с нормированным многочленом Чебышёва $\overline{T}_{n+1}^{[a; b]}$.

\textbf{Узлы Чебышёва} (корни $\overline{T}_{n+1}^{[a; b]}$):
\[x_m = \frac{a + b}{2} + \frac{b - a}{2} \cos\frac{\pi(2m - 1)}{2(n+1)}, \quad m = 1, \ldots, n+1\]

При таком выборе $Q_{n+1} = \overline{T}_{n+1}^{[a; b]}$ и достигается наилучшая оценка:
\[\|f - P_n\| \leq \frac{M_{n+1} (b-a)^{n+1} 2^{1-2(n+1)}}{(n+1)!}\]

Для любых других сеток правая часть оценки будет выше.

\textbf{Пример.} Функция $f(t) = 6t^3 + 3t + 3\sin(t^3 + 2)$ на $[-1; 1]$.

Сравниваются два приближения 4-й степени:
\begin{enumerate}
    \item $P_4$ --- равномерные узлы с шагом $h = 0{,}5$: $x_0 = -1, \ldots, x_4 = 1$.
    \item $P_{T4}$ --- узлы Чебышёва (корни $T_5$): $x_m = \cos\dfrac{\pi(2m-1)}{10}$, $m = 1, \ldots, 5$.
\end{enumerate}

Графики погрешностей $R_4(x) = |f(x) - P_4(x)|$ и $R_{T4}(x) = |f(x) - P_{T4}(x)|$ показывают, что наибольшая погрешность меньше для узлов Чебышёва, хотя в некоторых частях отрезка равномерные узлы дают лучшее приближение.

\hrule\bigskip

\subsection{10. Локальная интерполяция. Сплайны. Интерполяционный сплайн. Построение кубического сплайна. Погрешность интерполяции кубическим сплайном.}

\subsubsection{Локальная интерполяция}

Глобальная интерполяция (один многочлен на всю таблицу) имеет недостатки:
\begin{itemize}
    \item Сложное решение (высокая степень многочлена).
    \item Значительные погрешности (феномен Рунге).
\end{itemize}

\textbf{Локальная интерполяция} --- построение нескольких функций по участкам таблицы, из которых ``склеивается'' результирующая функция. Наиболее популярный метод --- \textbf{сплайн-интерполяция}.

\subsubsection{Интерполяционный сплайн}

\textbf{Сплайн} --- кусочно-полиномиальная функция. На каждом отрезке разбиения области совпадает с некоторым полиномом степени $m$, непрерывен со всеми производными до порядка $p$ включительно. Число $m - p$ --- \textbf{дефект сплайна}.

\textbf{Интерполяционный сплайн $S_m$ степени $m$:}
\begin{enumerate}
    \item На каждом $[x_{i-1}; x_i]$ совпадает с многочленом $P_{m,i}$ степени $m$.
    \item $S_m(x_i) = y_i$ (условия интерполяции).
    \item $S_m$ непрерывна на $[a; b]$ с производными до порядка $p$ включительно.
\end{enumerate}

\subsubsection{Построение кубического сплайна ($m = 3$, дефект 1)}

Шаги: $h_i = x_i - x_{i-1}$. На каждом $[x_{i-1}; x_i]$ сплайн определяется \textbf{формулой Эрмита}:
\[S_3(x) = y_{i-1} \frac{(x - x_i)^2(2(x - x_{i-1}) + h_i)}{h_i^3} + s_{i-1} \frac{(x - x_i)^2(x - x_{i-1})}{h_i^2} + y_i \frac{(x - x_{i-1})^2(2(x_i - x) + h_i)}{h_i^3} + s_i \frac{(x - x_{i-1})^2(x - x_i)}{h_i^2}\]

Параметры $s_i = S_3'(x_i)$ --- \textbf{наклоны} сплайна.

Формула Эрмита обеспечивает непрерывность $S_3$ и $S_3'$. Непрерывность $S_3''$ достигается подбором наклонов из условий $S_3''(x_i - 0) = S_3''(x_i + 0)$:
\[\frac{1}{h_i} s_{i-1} + 2\left(\frac{1}{h_i} + \frac{1}{h_{i+1}}\right) s_i + \frac{1}{h_{i+1}} s_{i+1} = 3\left(\frac{y_i - y_{i-1}}{h_i^2} + \frac{y_{i+1} - y_i}{h_{i+1}^2}\right)\]

Получена система из $n-1$ уравнений для $n+1$ неизвестных ($s_0, \ldots, s_n$). Для замыкания нужны \textbf{два граничных условия}.

\textbf{Виды граничных условий:}
\begin{enumerate}
    \item \textbf{На производные:} $s_0 = f'(x_0)$, $s_n = f'(x_n)$.
    \item \textbf{Свободного провисания (естественный сплайн):} $S_3''(x_0) = S_3''(x_n) = 0$:
    \[2s_0 + s_1 = 3\frac{y_1 - y_0}{h_1}, \quad s_{n-1} + 2s_n = 3\frac{y_n - y_{n-1}}{h_n}\]
    \item \textbf{``Отсутствия узла'':} $P_{3,1} \equiv P_{3,2}$ и $P_{3,n-1} \equiv P_{3,n}$ (непрерывность 3-й производной в $x_1$ и $x_{n-1}$).
    \item \textbf{Периодические} ($y_0 = y_n$): $s_0 = s_n$ и условие из $S_3''(x_0) = S_3''(x_n)$.
\end{enumerate}

Добавляя два граничных уравнения, получаем систему $As = b$ с трёхдиагональной (или почти трёхдиагональной) матрицей.

\subsubsection{Погрешность интерполяции кубическим сплайном}

\textbf{Теорема:} пусть $f$ имеет непрерывные производные до 4-го порядка на $[a; b]$. Тогда для кубического сплайна (кроме условий свободного провисания):
\[\Delta(S_3) = \max_{x \in [a; b]} |f(x) - S_3(x)| \leq C_4 M_4 h_{\max}^4\]
\[\Delta(S_3^{(k)}) = \max_{x \in [a; b]} |f^{(k)}(x) - S_3^{(k)}(x)| \leq C_k M_4 h_{\max}^{4-k}, \quad k = 1, 2, 3\]

где $M_4 = \max |f^{(4)}(x)|$, $h_{\max} = \max h_i$, $C_k > 0$ --- константы.

Сплайны приближают не только саму функцию, но и её производные до третьего порядка.

\hrule\bigskip

\subsection{11. Квадратичные сплайны. Примеры построения.}

\subsubsection{Квадратичные сплайны (стыковка в узлах)}

Сплайны второй степени: на каждом $[x_{i-1}; x_i]$ --- квадратный трёхчлен $P_{2,i}(x) = a_i x^2 + b_i x + c_i$.

Для $n$ отрезков --- $n$ трёхчленов, $3n$ неизвестных.

\textbf{Условия интерполяции} ($2n$ уравнений):
\[P_{2,1}(x_0) = y_0, \quad P_{2,1}(x_1) = y_1, \quad P_{2,2}(x_1) = y_1, \quad \ldots, \quad P_{2,n}(x_n) = y_n\]

\textbf{Непрерывность производной} ($n-1$ уравнений, дефект 1):
\[P_{2,i}'(x_i) = P_{2,i+1}'(x_i), \quad i = 1, \ldots, n-1\]

Итого $3n - 1$ уравнений. Не хватает одного --- добавляется \textbf{граничное условие}, например: $P_{2,1}'(x_0) = y_0'$.

\subsubsection{Квадратичные сплайны (стыковка в промежуточных точках)}

Вводятся \textbf{точки соединения} --- середины отрезков:
\[\tilde{x}_i = \frac{x_{i-1} + x_i}{2}, \quad i = 1, \ldots, n, \quad \tilde{x}_0 = x_0, \quad \tilde{x}_{n+1} = x_n\]

``Кусочки'' сплайна --- $n+1$ квадратный трёхчлен $P_{2,i}$ на отрезках $[\tilde{x}_{i-1}; \tilde{x}_i]$. Неизвестных: $3n + 3$.

\textbf{Уравнения:}

\begin{enumerate}
    \item \textbf{Интерполяция} ($n+1$ уравнений):
    \[P_{2,1}(x_0) = y_0, \quad P_{2,2}(x_1) = y_1, \quad \ldots, \quad P_{2,n+1}(x_n) = y_n\]

    \item \textbf{Непрерывность сплайна} ($n$ уравнений):
    \[P_{2,i}(\tilde{x}_i) = P_{2,i+1}(\tilde{x}_i), \quad i = 1, \ldots, n\]

    \item \textbf{Непрерывность производной} ($n$ уравнений, дефект 1):
    \[P_{2,i}'(\tilde{x}_i) = P_{2,i+1}'(\tilde{x}_i), \quad i = 1, \ldots, n\]

    \item \textbf{Граничные условия} (2 уравнения), например:
    \[P_{2,1}'(x_0) = y_0', \quad P_{2,n+1}'(x_n) = y_n'\]
\end{enumerate}

Итого: $(n+1) + n + n + 2 = 3n + 3$ уравнений. Решая систему, находят коэффициенты всех квадратных трёхчленов.

\end{document}
